\documentclass[12pt,a4paper]{article}

% --- PACKAGES DE BASE ---
\usepackage[utf8]{inputenc}
\usepackage[T1]{fontenc}
\usepackage[french]{babel}
\usepackage{geometry}
\geometry{top=2.5cm, bottom=2.5cm, left=2.5cm, right=2.5cm}
\usepackage{graphicx}
\usepackage{booktabs}
\usepackage{xcolor}
\usepackage{hyperref}
\usepackage{listings}
\usepackage{fancyhdr}
\usepackage{amsmath}
\usepackage{enumitem}

% --- CONFIGURATION GRAPHIQUE & COULEURS ---
\definecolor{primary}{HTML}{1D4ED8}     % Bleu primaire (Tailwind blue-700)
\definecolor{secondary}{HTML}{0A2647}   % Bleu nuit (Sidebar)
\definecolor{darkslate}{HTML}{1E293B}   % Slate 800
\definecolor{lightbg}{HTML}{F8FAFC}     % Slate 50
\definecolor{bordercolor}{HTML}{E2E8F0} % Slate 200
\definecolor{accentcolor}{HTML}{059669} % Vert émeraude

\hypersetup{
    colorlinks=true,
    linkcolor=primary,
    filecolor=primary,      
    urlcolor=primary,
    pdftitle={IncidentFlow - Suggestions d'Ameliorations},
    pdfpagemode=FullScreen,
}

% --- EN-TÊTES ET PIEDS DE PAGE ---
\pagestyle{fancy}
\fancyhf{}
\lhead{\textcolor{secondary}{\textbf{IncidentFlow}}}
\rhead{Suggestions d'Améliorations}
\cfoot{\thepage}
\renewcommand{\headrulewidth}{0.4pt}

% --- INFORMATIONS DU DOCUMENT ---
\title{
    \vspace{1cm}
    \textbf{\Huge IncidentFlow}\\[0.5cm]
    \Large Plateforme Dynamique de Gestion d'Incidents\\[1.5cm]
    \textbf{\Large Catalogue des Suggestions d'Améliorations et Perspectives d'Évolution}
    \vspace{1.5cm}
}
\author{\textbf{Anas HADDOU} \\ Élève Ingénieur en Génie Informatique}
\date{\today}

\begin{document}

% Page de garde
\maketitle
\thispagestyle{empty}
\clearpage

% Table des matières
\tableofcontents
\clearpage

% --- SECTION 1 : INTRODUCTION ---
\section{Introduction}
Ce document regroupe un ensemble de suggestions, d'améliorations et de perspectives d'évolutions pour le projet \textbf{IncidentFlow}. Il sert de guide pour enrichir le rapport final de stage d'ingénieur et orienter les futurs développements de l'application. 

L'architecture actuelle de l'application, basée sur un backend \textbf{Spring Boot}, un frontend \textbf{React (Vite)}, une persistance relationnelle \textbf{PostgreSQL} et une gestion d'identités \textbf{Keycloak}, offre une extensibilité naturelle pour intégrer ces nouveautés sans casser les processus métiers déjà implémentés.

---

% --- SECTION 2 : AMÉLIORATIONS FONCTIONNELLES & MÉTIER ---
\section{Améliorations Fonctionnelles et Métier}
Ces propositions visent à enrichir la logique métier d'IncidentFlow pour répondre aux exigences réelles des équipes d'exploitation ITIL.

\begin{description}
    \item[Gestion des SLA (Service Level Agreement) et Compte à Rebours :] \hfill \\
    Intégrer la notion de délai d'intervention et de résolution maximal selon la catégorie de l'incident et sa priorité.
    \begin{itemize}
        \item \textit{Mise en œuvre :} Stockage en base de données d'une heure limite calculée automatiquement à la création (ex: \texttt{temps\_creation + 2h} pour priorité Critique).
        \item \textit{Impact UX :} Affichage d'un minuteur ou d'une barre de progression colorée (vert $\rightarrow$ orange $\rightarrow$ rouge clignotant) sur les fiches des incidents pour alerter les opérateurs visuellement avant le dépassement du délai (SLA Breach).
    \end{itemize}

    \item[Moteur d'Escalade Automatique :] \hfill \\
    Si le SLA menace d'être dépassé sans qu'aucune action n'ait été enregistrée sur le ticket.
    \begin{itemize}
        \item \textit{Mise en œuvre :} Tâche planifiée en arrière-plan (\textit{Spring Scheduler}) qui vérifie périodiquement les incidents non résolus.
        \item \textit{Comportement :} Réaffectation automatique du ticket à un niveau de support supérieur (ex: Support N1 $\rightarrow$ Support N2) ou envoi d'une alerte prioritaire par e-mail au Responsable.
    \end{itemize}

    \item[Moteur de Validation Formelle des Workflows :] \hfill \\
    Permettre aux administrateurs de créer des processus dynamiques en limitant les risques de mauvaise configuration.
    \begin{itemize}
        \item \textit{Mise en œuvre :} Intégration d'un algorithme de parcours de graphe (BFS/DFS) lors de la sauvegarde d'un workflow dans le backend.
        \item \textit{Objectif :} Détecter les anomalies structurelles : états orphelins (impossibles à atteindre), puits de blocage (états sans aucune transition de sortie autre que vers eux-mêmes) et absence de chemin vers l'état final (clôture).
    \end{itemize}

    \item[Classification et Tri Intelligent par IA (Perspectives d'Évolution) :] \hfill \\
    Bien que l'Intelligence Artificielle ait été classée hors périmètre de la V1, elle constitue une piste d'évolution majeure.
    \begin{itemize}
        \item \textit{Mise en œuvre :} Intégration d'une API d'IA locale (via Ollama avec un modèle Mistral/Llama) ou distante (Google Gemini API).
        \item \textit{Fonctionnalité :} Analyse sémantique du titre et de la description saisis par l'utilisateur pour suggérer automatiquement la catégorie appropriée, le niveau de priorité de l'incident et affecter le ticket au bon service de support.
    \end{itemize}

    \item[Versionnage et Historisation des Workflows (Audit Trail) :] \hfill \\
    Que se passe-t-il lorsqu'un administrateur modifie la structure d'un workflow alors que des incidents y sont rattachés ?
    \begin{itemize}
        \item \textit{Mise en œuvre :} Implémenter un système de versioning (ex: Workflow v1.0, Workflow v1.1).
        \item \textit{Fonctionnalité :} Les incidents en cours de traitement continuent sur la version du workflow avec laquelle ils ont été créés, tandis que les nouveaux incidents adoptent la dernière version publiée. Cela préserve l'intégrité de l'historique et évite les corruptions de transitions.
    \end{itemize}
\end{description}

---

% --- SECTION 3 : AMÉLIORATIONS FRONTEND & ERGONOMIQUES ---
\section{Améliorations Frontend et Ergonomiques (UX/UI)}
L'expérience utilisateur est primordiale pour garantir l'adoption de la plateforme par les équipes de support sous pression opérationnelle.

\begin{description}
    \item[Concepteur Visuel de Workflows (React Flow) :] \hfill \\
    Remplacer la création de workflows par formulaire par une interface visuelle interactive de type glisser-déposer (Drag-and-Drop).
    \begin{itemize}
        \item \textit{Technologie :} Bibliothèque \texttt{React Flow} (ou \texttt{GoJS}).
        \item \textit{Ergonomie :} L'administrateur dessine le workflow : chaque état est une boîte graphique (nœud), et chaque transition est représentée par une ligne reliant deux nœuds. Cela réduit drastiquement les erreurs de saisie et rend l'administration accessible aux utilisateurs métiers.
    \end{itemize}

    \item[Notifications Instantanées en Temps Réel (WebSockets / SSE) :] \hfill \\
    Recevoir des mises à jour dynamiques sans rechargement manuel de la page.
    \begin{itemize}
        \item \textit{Technologie :} \textit{Server-Sent Events} (SSE) pour une communication unidirectionnelle légère, ou \textit{WebSockets} via Spring Boot et SockJS.
        \item \textit{Cas d'usage :} Si un opérateur se voit attribuer un ticket ou si un commentaire est ajouté par un autre utilisateur, une notification push native s'affiche instantanément sur son écran.
    \end{itemize}

    \item[Barre de Commande Universelle (Command Palette) :] \hfill \\
    Inspirée des outils de développement modernes (VS Code, Slack).
    \begin{itemize}
        \item \textit{Mise en œuvre :} Combinaison clavier \texttt{Ctrl + K} ouvrant une fenêtre de recherche modale au centre de l'écran.
        \item \textit{Intérêt :} Permet de taper un mot-clé pour sauter immédiatement vers un incident (\texttt{\#INC-1243}), rechercher un compte utilisateur ou déclencher des actions rapides (ex: déclarer un incident) sans utiliser la souris.
    \end{itemize}

    \item[Mode Sombre (Dark Mode) et Accessibilité :] \hfill \\
    Les équipes de support travaillent souvent en horaires décalés (24/7).
    \begin{itemize}
        \item \textit{Mise en œuvre :} Intégration de variables CSS globales pour basculer facilement le thème.
        \item \textit{Conformité :} Respect des règles d'accessibilité WCAG (contraste des textes, navigation au clavier et lecture d'écran pour les agents malvoyants).
    \end{itemize}
\end{description}

---

% --- SECTION 4 : AMÉLIORATIONS TECHNIQUES, DEVOPS & PERFORMANCES ---
\section{Améliorations Techniques, DevOps et Performances}
Pour garantir la robustesse du système en production (scalabilité, tolérance aux pannes et sécurité).

\begin{description}
    \item[Mise en cache avec Redis :] \hfill \\
    Le workflow dynamique nécessite des calculs SQL fréquents pour valider les états à chaque chargement d'incident.
    \begin{itemize}
        \item \textit{Architecture :} Ajout d'un conteneur \texttt{incidentflow-redis} au fichier Docker Compose.
        \item \textit{Bénéfice :} Mise en cache des listes d'états et de transitions par catégorie. Le backend ne sollicite PostgreSQL que lors des mises à jour administratives du workflow, divisant le temps de réponse de l'API par 5.
    \end{itemize}

    \item[Synchronisation de Session par Webhooks Keycloak :] \hfill \\
    Si un compte utilisateur est désactivé ou si ses rôles changent dans l'administration Keycloak, l'application Spring Boot doit réagir instantanément.
    \begin{itemize}
        \item \textit{Mécanisme :} Configuration de webhooks Keycloak ou écouteurs d'événements SPI (Service Provider Interface).
        \item \textit{Résultat :} Le backend reçoit une alerte et invalide immédiatement le jeton JWT de l'utilisateur actif concerné, provoquant sa déconnexion forcée du frontend pour des raisons de sécurité.
    \end{itemize}

    \item[Tests d'Intégration Automatisés avec Testcontainers :] \hfill \\
    Garantir que les modifications de code ne cassent pas l'interaction avec la base de données ou la sécurité Keycloak.
    \begin{itemize}
        \item \textit{Technologie :} \texttt{Testcontainers} avec Spring Boot.
        \item \textit{Intérêt :} Lors des tests d'intégration (lancement de Maven en CI/CD), Testcontainers démarre de vrais conteneurs éphémères PostgreSQL et Keycloak, permettant de tester le code en situation réelle de déploiement.
    \end{itemize}

    \item[Observabilité Applicative (Spring Boot Actuator, Prometheus \& Grafana) :] \hfill \\
    Suivre l'utilisation et la santé des conteneurs en temps réel.
    \begin{itemize}
        \item \textit{Architecture :} Ajout d'agents de collecte de métriques.
        \item \textit{Bénéfices :} Création de graphiques de supervision pour analyser la consommation de RAM/CPU du backend, le nombre de requêtes HTTP par seconde, les erreurs 500 et les goulots d'étranglement de la base de données.
    \end{itemize}
\end{description}

---

% --- SECTION 5 : SYNTHÈSE POUR LE RAPPORT ---
\section{Synthèse pour le Rapport Final}
Dans le cadre de votre soutenance ou de la rédaction du mémoire de stage, il est recommandé de présenter ces propositions comme une section dédiée nommée \textbf{« Limites actuelles et Perspectives d'Évolution »}. 

Cela démontre aux examinateurs que :
\begin{enumerate}
    \item Vous avez conscience des contraintes d'échelle et des problématiques de production (sécurité avancée, performances SQL).
    \item Vous maîtrisez l'état de l'art technologique (SSE, Redis, Testcontainers, React Flow).
    \item Vous savez analyser de manière critique votre propre travail pour proposer des axes d'amélioration structurés.
\end{enumerate}

\end{document}
